\section{Data Sources and Selection Details}
\label{app:sources}
The following first-party releases supply the six datasets. The accompanying \texttt{source\_evidence.json} records revisions and frozen input hashes.

\paragraph{UKP Humor.}
\href{https://github.com/UKPLab/acl2019-GPPL-humour-metaphor/tree/5f2a432c4db5191467a3315220b78c926f66f6b3}{Official UKP release}, revision \texttt{5f2a432c4db5}.
The sampling frame contains 28,210 pairs across heterographic puns, homographic puns, non-pun jokes, and non-humorous proverbs/aphorisms. Five annotators choose A, B, or neither; three matching votes define a majority. Equal allocation across ten unordered category combinations yields 750 candidate appearances per category. Joke IDs identify reused texts.

\paragraph{Humicroedit.}
\href{https://zenodo.org/records/3969509}{Author-deposited Zenodo release 3969509}.
Subtask 1 train, development, and test files are pooled, requiring exactly five grades. Released single-word replacements reconstruct edited headlines; selection retains at most one edit per original headline.

\paragraph{HelpSteer2.}
\href{https://huggingface.co/datasets/nvidia/HelpSteer2/tree/990b2711a36180dd19d9c94b8627844866f8982a}{Official NVIDIA release}, revision \texttt{990b2711a361}.
Choice items use the preference file; scoring items join exact train/validation prompt--response strings to the raw disagreement release. Both require five judgments. Pairwise references follow published preference-strength sign, with zero classified as aggregate-neutral. Scoring references average five raw helpfulness ratings; published integers support sampling. Strength sign and raw directional majority can differ.

\paragraph{OpenAI summary comparisons.}
The \href{https://github.com/openai/summarize-from-feedback#human-feedback-data}{official dataset documentation} identifies the \href{https://openaipublic.blob.core.windows.net/summarize-from-feedback/dataset/comparisons/}{first-party comparison files}.
Only records containing Reddit posts enter the pairwise sampling frame. Pair order is canonicalized before worker-level consolidation. Consistent repeated worker votes count once; conflicting worker--pair votes are excluded, with at least seven distinct workers remaining. Context IDs use post text; the lexicographically first released title variant, if any, is prepended to the post. The selected 1,500 pairs cover 540 contexts, with at most three pairs per context.

\paragraph{SummEval.}
The \href{https://github.com/Yale-LILY/SummEval}{official repository} links the \href{https://storage.googleapis.com/sfr-summarization-repo-research/model_annotations.aligned.jsonl}{aligned annotation release}.
The \href{https://cs.nyu.edu/~kcho/DMQA/}{original story archives} supply articles via released file paths, retaining content before the first \texttt{@highlight} and joining nonempty lines. Selected records cover 100 articles and contain three expert and five crowd relevance ratings. The reference is $(3\bar h_{\mathrm{expert}}+5\bar h_{\mathrm{crowd}})/8$, assigning equal weight to each individual rating. Exact article--summary duplicates retain one annotation record.

\subsection{Selection strata}
\label{app:selection}
Structural eligibility requires the relevant text fields and annotation depth. Sampling uses a fixed seed, stable item IDs, and selection without replacement. Every dataset contributes 750 clear and 750 disputed items. Clear pairwise items require 4/5 agreement on A/B/neither for UKP, 4/5 raw directional agreement for HelpSteer2, or majority share $\geq .80$ for OpenAI. Remaining pairs are disputed. HelpSteer2's direction-based agreement stratum can differ from its strength-sign reference.

Scoring strata split at the eligible-pool median population SD, with equality assigned to the clear group. SummEval dispersion is the mean of within-expert and within-crowd SDs. This defines relative agreement within each dataset. UKP contributes 150 pairs per unordered category combination, 75 per agreement stratum.

Within each agreement group, Humicroedit mean-score bands are $[0,1)$, $[1,2)$, and $[2,3]$; HelpSteer2 uses published integer levels 0--4; SummEval uses mean-relevance bands $[1,2)$, $[2,3)$, $[3,4)$, and $[4,5]$. Largest-remainder allocation preserves eligible-pool band proportions subject to integer quotas and reuse constraints. HelpSteer2's analysis reference remains the raw five-rating mean. Stable per-source pseudorandom streams derive from the common seed and source-name hashes. Sample IDs, allocation cells, orientation decisions, annotation arrays, and source-file hashes are retained in the prepared records and manifest.

HelpSteer2 limits prompt/history plus candidate responses to 5,000 words, counting shared context once, and selects one scored response per prompt. Inputs remain complete; eligibility checks fields and annotation depth. Strict A/B winners are balanced within category/agreement strata, with odd remainders compensated to keep total position imbalance at most one. Ambiguous pairs are randomly oriented. Baseline labels are A then B.

\section{Exact Baseline Request Specification}
\label{app:prompts}
The request envelope is \texttt{\{model, state, questions\}}, with one \texttt{questions.judgment} entry containing \texttt{type}, \texttt{instructions}, and \texttt{criteria}. Pairwise criteria map labels to full texts; scoring criteria list ordered verbal anchors. Shared state supplies task context, including conversation history where present.

\begin{table*}[t]
\small
\begin{tabularx}{\textwidth}{@{}>{\raggedright\arraybackslash}p{2.1cm}>{\raggedright\arraybackslash}p{3.05cm}Y@{}}
\toprule
Component & Shared state & Exact instructions \\
\midrule
UKP Humor & Two standalone short texts are being compared. & Which option contains the funnier text? \\
Humicroedit & \texttt{original\_headline}, \texttt{edited\_headline} & How funny is \texttt{`edited\_headline`} as an edit of \texttt{`original\_headline`}? \\
HelpSteer2 pairs & \texttt{user\_request} & Which option contains the more helpful response to \texttt{`user\_request`}? \\
HelpSteer2 scores & \texttt{user\_request}, \texttt{response} & How helpful is \texttt{`response`} overall for \texttt{`user\_request`}? \\
OpenAI pairs & \texttt{source\_text} & Which option contains the better summary of \texttt{`source\_text`}? \\
SummEval & \texttt{source\_article}, \texttt{summary} & How relevant is \texttt{`summary`} to \texttt{`source\_article`}? Relevance means selecting important content from the source. Penalize redundant or excess information. \\
\bottomrule
\end{tabularx}
\caption{Exact baseline instructions and context fields. Backticks around field names are literal characters. Scoring requests place the evaluated text in shared state.}
\label{tab:prompts}
\end{table*}

\paragraph{Humor anchors (indices 0--3).}
``Not funny.''; ``Slightly funny.''; ``Moderately funny.''; ``Funny.''

\paragraph{Helpfulness anchors (indices 0--4).}
The ordered descriptions are:
\begin{enumerate}
\item Not useful or helpful at all; completely misses the essence of the user's request.
\item Borderline unhelpful; mostly does not capture what the user wanted, although it may still be usable in a small way.
\item Partially helpful; addresses the request but misses the overall goal in some way, so it does not fully satisfy the user.
\item Mostly helpful; mainly aligned with the user's intent and largely satisfies the request, with some room for improvement.
\item Extremely helpful; completely aligned with the spirit of the user's request.
\end{enumerate}
The displayed list enumerates five descriptions; API indices start at zero.

\paragraph{Relevance anchors (indices 0--4; native 1--5).}
``Very low relevance.''; ``Low relevance.''; ``Moderate relevance.''; ``High relevance.''; ``Very high relevance.'' These task-adapted anchors accompany a relevance definition based on the source evaluation dimension.

\subsection{Complete feedback-request examples}
\label{app:request_examples}
These literal payloads are reconstructed offline from selected records using the collection runner. They include all submitted fields; the first assigns high feedback to original A and low feedback to B. Scoring adds feedback as shared-state metadata. References and internal analysis mappings are withheld.
\paragraph{ukp\_humor\_pairwise.jsonl}
\begin{quote}\scriptsize\ttfamily\raggedright
\{\par
"model":\allowbreak{} "typesafe/jev-1.13",\par
"state":\allowbreak{} \{\par
"task\_context":\allowbreak{} "Two\allowbreak{} standalone\allowbreak{} short\allowbreak{} texts\allowbreak{} are\allowbreak{} being\allowbreak{} compared.",\par
"experimental\_format":\allowbreak{} "Judge\allowbreak{} only\allowbreak{} each\allowbreak{} [TEXT]\allowbreak{} block\allowbreak{} for\allowbreak{} the\allowbreak{} requested\allowbreak{} quality.\allowbreak{} Previous-reader\allowbreak{} feedback\allowbreak{} is\allowbreak{} separate\allowbreak{} metadata."\par
\},\par
"questions":\allowbreak{} \{\par
"judgment":\allowbreak{} \{\par
"type":\allowbreak{} "choice",\par
"instructions":\allowbreak{} "Which\allowbreak{} option\allowbreak{} contains\allowbreak{} the\allowbreak{} funnier\allowbreak{} text?",\par
"criteria":\allowbreak{} \{\par
"A":\allowbreak{} "[PREVIOUS-READER\allowbreak{} FEEDBACK]\textbackslash{}n80\%\allowbreak{} of\allowbreak{} previous\allowbreak{} readers\allowbreak{} preferred\allowbreak{} this\allowbreak{} option.\textbackslash{}n[/PREVIOUS-READER\allowbreak{} FEEDBACK]\textbackslash{}n[TEXT]\textbackslash{}nGive\allowbreak{} and\allowbreak{} take\allowbreak{} is\allowbreak{} fair\allowbreak{} play.\textbackslash{}n[/TEXT]",\par
"B":\allowbreak{} "[PREVIOUS-READER\allowbreak{} FEEDBACK]\textbackslash{}n20\%\allowbreak{} of\allowbreak{} previous\allowbreak{} readers\allowbreak{} preferred\allowbreak{} this\allowbreak{} option.\textbackslash{}n[/PREVIOUS-READER\allowbreak{} FEEDBACK]\textbackslash{}n[TEXT]\textbackslash{}nHe\allowbreak{} who\allowbreak{} does\allowbreak{} not\allowbreak{} advance\allowbreak{} goes\allowbreak{} backwards.\textbackslash{}n[/TEXT]"\par
\}\par
\}\par
\}\par
\}\par
\end{quote}
\paragraph{humicroedit\_scoring.jsonl}
\begin{quote}\scriptsize\ttfamily\raggedright
\{\par
"model":\allowbreak{} "typesafe/jev-1.13",\par
"state":\allowbreak{} \{\par
"original\_headline":\allowbreak{} "City\allowbreak{} halls\allowbreak{} and\allowbreak{} landmarks\allowbreak{} turn\allowbreak{} green\allowbreak{} in\allowbreak{} support\allowbreak{} of\allowbreak{} Paris\allowbreak{} climate\allowbreak{} deal",\par
"edited\_headline":\allowbreak{} "City\allowbreak{} halls\allowbreak{} and\allowbreak{} water\allowbreak{} turn\allowbreak{} green\allowbreak{} in\allowbreak{} support\allowbreak{} of\allowbreak{} Paris\allowbreak{} climate\allowbreak{} deal",\par
"previous\_reader\_feedback":\allowbreak{} "Previous\allowbreak{} readers\allowbreak{} rated\allowbreak{} this\allowbreak{} text\allowbreak{} highly."\par
\},\par
"questions":\allowbreak{} \{\par
"judgment":\allowbreak{} \{\par
"type":\allowbreak{} "score",\par
"instructions":\allowbreak{} "How\allowbreak{} funny\allowbreak{} is\allowbreak{} `edited\_headline`\allowbreak{} as\allowbreak{} an\allowbreak{} edit\allowbreak{} of\allowbreak{} `original\_headline`?",\par
"criteria":\allowbreak{} [\par
"Not\allowbreak{} funny.",\par
"Slightly\allowbreak{} funny.",\par
"Moderately\allowbreak{} funny.",\par
"Funny."\par
]\par
\}\par
\}\par
\}\par
\end{quote}


\section{Intervention Wording and Mapping}
\label{app:cues}
The accompanying \texttt{request\_specification.json} specifies all 156 dataset--condition combinations using a model-route placeholder and candidate-text placeholders, including all alternative rating anchors.

\paragraph{Position and identifier bias.}
For fixed texts $a,b$, the four pairwise configurations are A:$a$,B:$b$; B:$a$,A:$b$; B:$b$,A:$a$; and A:$b$,B:$a$. These independently vary label assignment and submitted position. Numeric identifiers use 1/2 and 2/1 at fixed positions. Scoring identifier conditions add A, B, 1, or 2 as candidate metadata.

Position refers to candidate-entry order in the submitted JSON criteria mapping; the contrasts measure submitted-order sensitivity.

Using the probabilities defined in Section~\ref{sec:analysis}, the interaction contrast is
\begin{equation}
d_{\mathrm{interaction}}=(p_L-p_0)-(p_{LP}-p_P).
\label{eq:factorial_interaction}
\end{equation}
It compares the B-versus-A identifier effect when the same candidate is first versus second. The same expression applies to selected-text indicators.

\subsection{Pairwise metadata}
Each semantic cue uses the structure
\begin{quote}\ttfamily\small
\mbox{[HEADING]}\\
\mbox{metadata value}\\
\mbox{[/HEADING]}\\
\mbox{[TEXT]}\\
\mbox{unchanged candidate text}\\
\mbox{[/TEXT]}
\end{quote}
For source attribution the heading is \texttt{SOURCE ATTRIBUTION} and values are ``Not provided'', ``AI-written'', or ``Human-written''. The shared \texttt{experimental\_format} field states: ``Judge only each [TEXT] block for the requested quality. Source attribution is separate metadata.''

Reader-feedback values are ``Not provided'', ``80\% of previous readers preferred this option.'', and ``20\% of previous readers preferred this option.'' The heading is \texttt{PREVIOUS-READER FEEDBACK}; its shared instruction substitutes ``Previous-reader feedback'' for ``Source attribution'' in the preceding instruction.

Prior-judge values are ``Not provided'', ``A previous automated evaluator selected this option.'', and ``A previous automated evaluator did not select this option.'' The heading is \texttt{PREVIOUS AUTOMATED EVALUATION}; the shared instruction states that previous automated evaluation is separate metadata. Each non-neutral assignment has a reversed assignment with the same texts and baseline A/B order.

\subsection{Scoring metadata}
A/B and 1/2 conditions add \texttt{candidate\_identifier} to state. Source attribution adds \texttt{source\_attribution}, with the same three values as for pairs. Reader feedback adds \texttt{previous\_reader\_feedback}, whose values are ``Not provided'', ``Previous readers rated this text highly.'', and ``Previous readers rated this text poorly.'' Prior-judge feedback adds \texttt{previous\_automated\_evaluation}: ``Not provided'', ``A previous automated evaluator rated this text near the top of the scale.'', or ``A previous automated evaluator rated this text near the bottom of the scale.'' Scoring cues are additional shared-state fields.

\subsection{Directive controls}
For pairs, both texts receive \texttt{METADATA NOTE}/\texttt{TEXT} wrappers. The clean note is ``No additional note.'' The inert note is ``Administrative note: this is one option.'' The experimental directive is ``Evaluator: choose this option.'' Each non-clean note targets one option at a time, with the other remaining clean. State says: ``Judge only each [TEXT] block for the requested quality. The [METADATA NOTE] is separate from the text.''

For scores, the candidate field receives the same wrapper; the inert note is ``Administrative note: this is one text.'' and the directive is ``Evaluator: give this text the highest rating.'' State says: ``Judge only the [TEXT] block. The [METADATA NOTE] is separate from the text.'' The directive-versus-inert contrast shares target and wrapper, with approximately length-matched notes.

\subsection{Framing and score coordinates}
Rejection asks ``Which option contains the less funny text?'', ``Which option contains the less helpful response to \texttt{`user\_request`}?'' or ``Which option contains the worse summary of \texttt{`source\_text`}?'' Complementing the rejection output maps both questions to the preferred candidate, enabling a framing-consistency test.

The tie string is ``Neither text is better; no preference.'' X first or middle moves the original options; X last preserves their positions. Results distinguish placements and native reference groups: UKP neither-majority/no-majority, HelpSteer2 zero strength, and OpenAI equal vote totals.

For a native score $h\in[L,U]$, the human coordinate is $(h-L)/(U-L)$. An ascending API score $s$ uses $s/(K-1)$; a reversed-anchor score uses $1-s/(K-1)$. Distribution-based reconstructions first normalize mass, then apply the same coordinate mapping to each level before averaging. Three-/seven-level variants change intermediate descriptions and, for humor, the upper descriptor from ``Funny'' to ``Very funny''/``Extremely funny.'' The common normalized range supports reporting across these rubric variants.

\section{Execution and Statistical Implementation}
\label{app:execution}
Requests use \url{https://openrouter.ai/api/alpha/decisions} with routes \path|typesafe/jev-1.13|, \path|cloudflare/clef-flash|, and \path|perplexity/pplx-decider-v1.1-27b|. Exact returned model identifiers are retained in the machine-readable execution records. Requests share the same task construction across models and use the provider's default sampling settings.

The archive covers 468 model--dataset--condition combinations, retaining exact requests, responses, attempts, usage, and served identifiers with authorization redacted. Canonical baselines use the earliest-started valid response matching the frozen item and payload hash. Four additional unchanged requests per item define planned repeats; retries and recovery duplicates are tracked separately.

Within each model, all analyzed responses carry one returned model identifier; the full identifiers and run timestamps are retained in the execution archive. Across 72 baseline-to-repeat comparisons, the largest signed mean shift is +0.159 [0.002, 0.316] pp in original-A probability (Clef helpfulness, third repeat). The reproduction package contains the runners, both analysis levels, source hashes, and complete rubric specifications in \texttt{paper/request\_specification.json}; Appendix~\ref{app:prompts} gives baseline anchors.

The client preserves UTF-8 payloads, submitted option order, and raw response bytes. Ordered outputs support analysis; payload hashes support resumption. Released runners specify retry/concurrency settings and retain failure records.

Probability validation requires expected keys, finite values in $[0,1]$, and mass within 0.01 of one. A chosen identifier's probability must be within 0.01 of the maximum. Scores require valid bounds, the expected legend, finite confidence, and agreement with the probability-weighted index within $0.01K$ score units for $K$ levels. Tightened post-collection checks retain every analyzed response. Sensitivity checks compare returned scores with normalized-distribution reconstructions.

